% Template for a single lecture note.
% Copy this file into lectures/<category>/<slug>/notes.tex and fill it in.
% Compile with: latexmk -xelatex -interaction=nonstopmode notes.tex
\documentclass[11pt]{article}
\usepackage{../../../tex/lecturestyle}

\title{LECTURE TITLE}
\author{}
\date{}

\begin{document}
\thispagestyle{fancy}

% #1 category, #2 title, #3 difficulty 1-9
\lecturetitleband{CATEGORY NAME}{LECTURE TITLE}{N}

\section{Overview}
Short original overview of the problem this technique solves.

\section{Definitions}
\begin{definitionbox}
State the key definitions here.
\end{definitionbox}

\section{The Technique}
Explain the idea, invariants, and correctness argument in your own words.

\begin{theorembox}
State and (briefly) justify the main correctness/complexity claim.
\end{theorembox}

\section{Complexity}
Time and space complexity discussion.

% Illustrations: use TikZ inside a [H] figure, with the shared styles
% (lnode, lcell, lhi, lhi2, larrow, lthin, llabel, lvert) from lecturestyle.sty.
%
% \begin{figure}[H]
% \centering
% \begin{tikzpicture}
%   \node[lnode] (a) {start}; \node[lnode,lhi,right=of a] (b) {end};
%   \draw[larrow] (a) -- (b);
% \end{tikzpicture}
% \caption{What the picture shows and why it matters.}\label{fig:example}
% \end{figure}

% Code: example.rs / example.py must contain `PDFBEGIN` and `PDFEND` marker
% comments; only the lines between them are printed by \corelisting.
\section{Reference Implementations}

\codenote
\subsection{Rust}
\corelisting{Rust}{example.rs}

\subsection{Python}
\corelisting{Python}{example.py}

\section{Worked Example}
\begin{examplebox}
Walk through a small concrete input.
\end{examplebox}

\section{Further Reading}
\begin{itemize}[leftmargin=1.5em]
  \item AtCoder Algorithm Lectures index: \url{https://info.atcoder.jp/entry/algorithm_lectures/index}
\end{itemize}

\end{document}
